\chapter{2021年上海卷（春）}

\section{填空题}

\begin{problem}
等差数列 \(\{a_n\}\) 中，\(a_1=3\)，\(d=2\)，则 \(a_{10}=\fillinblank{}\).
\end{problem}

\begin{problem}
已知复数 \(z=1-3\i\)，则 \(|\bar z-\i|=\fillinblank{}\).
\end{problem}

\begin{problem}
不等式 \(\dfrac{2x+5}{x-2}<1\) 的解集为\fillinblank{}.
\end{problem}

\begin{problem}
已知圆柱的底面半径为 \(1\)，母线长为 \(2\)，则其侧面积为\fillinblank{}.
\end{problem}

\begin{problem}
直线 \(x=-2\) 与直线 \(\sqrt3x-y+1=0\) 的夹角为\fillinblank{}.
\end{problem}

\begin{problem}
方程组 \(\begin{cases}
a_1x+b_1y=c_1,\\
a_2x+b_2y=c_2
\end{cases}\) 无解，则 \(\begin{vmatrix}a_1&b_1\\a_2&b_2\end{vmatrix}=\fillinblank{}\).
\end{problem}

\begin{problem}
\((x+1)^n\) 的二项展开式中有且仅有 \(x^3\) 的系数为最大值，则 \(x^3\) 的系数为\fillinblank{}.
\end{problem}

\begin{problem}
已知函数 \(f(x)=3^x+\dfrac{a}{3^x+1}(a>0)\) 的最小值为 \(5\)，则 \(a=\fillinblank{}\).
\end{problem}

\begin{problem}
在无穷等比数列 \(\{a_n\}\) 中，\(\lim\limits_{n\to\infty}(a_1-a_n)=4\)，则 \(a_2\) 的取值范围为\fillinblank{}.
\end{problem}

\begin{problem}
某人某天需要运动总时长大于等于 \(60\) 分钟，现有五项运动可选，时间分别为 \(30\)，\(20\)，\(40\)，\(30\)，\(30\) 分钟，则有\fillinblank{} 种运动组合.
\end{problem}

\begin{problem}
已知椭圆 \(x^2+\dfrac{y^2}{b^2}=1(0<b<1)\) 的左，右焦点为 \(F_1\)，\(F_2\)，以 \(O\) 为顶点，\(F_2\) 为焦点作抛物线交椭圆于 \(P\)，且 \(\angle PF_1F_2=45^\circ\)，则抛物线的准线方程是\fillinblank{}.
\end{problem}

\begin{problem}
已知 \(\theta>0\)，对任意 \(n\in\mathbb N^*\)，总存在实数 \(\varphi\)，使得 \(\cos(n\theta+\varphi)<\dfrac{\sqrt3}{2}\)，则 \(\theta\) 的最小值是\fillinblank{}.
\end{problem}

\section{单选题}

\begin{problem}
下列函数中，在定义域内存在反函数的是（\quad）.
\choices{\(x^2\)}{\(\sin x\)}{\(2^x\)}{\(y=1\)}
\end{problem}

\begin{problem}
已知集合 \(A=\{x\mid x^2-x-2\ge0\}\)，\(B=\{x\mid x>-1\}\)，则（\quad）.
\choices{\(A\subseteq B\)}{\(\complement_{\mathbb R}A\subseteq\complement_{\mathbb R}B\)}{\(A\cap B=\varnothing\)}{\(A\cup B=\mathbb R\)}
\end{problem}

\begin{problem}
已知函数 \(y=f(x)\) 的定义域为 \(\mathbb R\)，下列是 \(f(x)\) 无最大值的充分条件是（\quad）.
\choices
    {\(f(x)\) 为偶函数且关于直线 \(x=1\) 对称}
    {\(f(x)\) 为偶函数且关于点 \((1,1)\) 对称}
    {\(f(x)\) 为奇函数且关于直线 \(x=1\) 对称}
    {\(f(x)\) 为奇函数且关于点 \((1,1)\) 对称}
\end{problem}

\begin{problem}
在 \(\triangle ABC\) 中，\(D\) 为 \(BC\) 中点，\(E\) 为 \(AD\) 中点. 判断：
\begin{circlelist}
\item 存在 \(\triangle ABC\)，使得 \(\overrightarrow{AB}\cdot\overrightarrow{CE}=0\)；
\item 存在 \(\triangle ABC\)，使得 \(\overrightarrow{CE}\parallel(\overrightarrow{CB}+\overrightarrow{CA})\).
\end{circlelist}
则（\quad）.
\choices{\(\circled{1}\) 成立，\(\circled{2}\) 成立}{\(\circled{1}\) 成立，\(\circled{2}\) 不成立}{\(\circled{1}\) 不成立，\(\circled{2}\) 成立}{\(\circled{1}\) 不成立，\(\circled{2}\) 不成立}
\end{problem}

\section{解答题}

\begin{problem}
四棱锥 \(P-ABCD\)，底面为正方形 \(ABCD\)，边长为 \(4\)，\(E\) 为 \(AB\) 中点，\(PE\perp\) 平面 \(ABCD\).
\begin{enumerate}
\item 若 \(\triangle PAB\) 为等边三角形，求四棱锥 \(P-ABCD\) 的体积；
\item 若 \(CD\) 的中点为 \(F\)，\(PF\) 与平面 \(ABCD\) 所成角为 \(45^\circ\)，求 \(PD\) 与 \(AC\) 所成角的大小.
\end{enumerate}
\end{problem}

\begin{problem}
已知 \(A\)，\(B\)，\(C\) 为 \(\triangle ABC\) 的三个内角，\(a\)，\(b\)，\(c\) 是其三条边，\(a=2\)，\(\cos C=-\dfrac14\).
\begin{enumerate}
\item 若 \(\sin A=2\sin B\)，求 \(b\)，\(c\)；
\item 若 \(\cos\left(A-\dfrac{\pi}{4}\right)=\dfrac45\)，求 \(c\).
\end{enumerate}
\end{problem}

\begin{problem}
团队在 \(O\) 点西侧，东侧 \(20\) 千米处设有 \(A\)，\(B\) 两站点，测量一点 \(P\) 满足 \(|PA|-|PB|=20\) 千米. 以 \(O\) 为原点，东侧为 \(x\) 轴正半轴，北侧为 \(y\) 轴正半轴建立平面直角坐标系，\(P\) 在北偏东 \(60^\circ\) 处.
\begin{enumerate}
\item 求双曲线标准方程和 \(P\) 点坐标；
\item 又在南侧，北侧 \(15\) 千米处设 \(C\)，\(D\) 两站点，若 \(|QA|-|QB|=30\) 千米，\(|QC|-|QD|=10\) 千米，求 \(|OQ|\) 和 \(Q\) 点位置.
\end{enumerate}
\end{problem}

\begin{problem}
已知函数 \(f(x)=|x+a|-a-x\).
\begin{enumerate}
\item 若 \(a=1\)，求函数定义域；
\item 若 \(a\ne0\)，且 \(f(ax)=a\) 有 \(2\) 个不同实数根，求 \(a\) 的取值范围；
\item 是否存在实数 \(a\)，使得函数 \(f(x)\) 在定义域内具有单调性？若存在，求出 \(a\) 的取值范围.
\end{enumerate}
\end{problem}

\begin{problem}
已知数列 \(\{a_n\}\) 满足 \(a_n\ge0\)，对任意 \(n\ge2\)，\(a_n\) 和 \(a_{n+1}\) 中存在一项使其为另一项与 \(a_{n-1}\) 的等差中项.
\begin{enumerate}
\item 已知 \(a_1=5\)，\(a_2=3\)，\(a_4=2\)，求 \(a_3\) 的所有可能取值；
\item 已知 \(a_1=a_4=a_7=0\)，\(a_2\)，\(a_5\)，\(a_8\) 为正数，证明：\(a_2\)，\(a_5\)，\(a_8\) 成等比数列，并求公比 \(q\)；
\item 已知数列中恰有 \(3\) 项为 \(0\)，即 \(a_r=a_s=a_t=0\)，\(2<r<s<t\)，且 \(a_1=1\)，\(a_2=2\)，求 \(a_{r+1}+a_{s+1}+a_{t+1}\) 的最大值.
\end{enumerate}
\end{problem}

